\section{Limitations and future directions}
\label{sec:limitations}

\paragraph{Closed-set evaluation.}
The headline 98.7\% and $r\!\approx\!1$ figures characterize classification
against genomes that are present in the reference. Because our test organisms
are drawn from the same catalogue used to build the training reference, these
numbers measure closed-set performance, not novel-genome generalization. We
showed the mechanism is learned composition rather than exact-match lookup
(Section~\ref{sec:lookup}), but the more demanding question---how accuracy
transfers to genomes of the same genera that were \emph{withheld} from
training---is the single most important open item. An out-of-genome
generalization experiment (holding out a subset of each genus's genomes from
training) is the natural next step and is required before the tokenization
advantage can be claimed for real-world discovery settings.

\paragraph{Simulated reads.}
Reads were generated with ART \citep{Huang2012}. Simulation may favour
alignment- and $k$-mer-based methods with structurally clean error models, and
does not capture all real sequencing artefacts. Our probe on a real mock
community (ZymoBIOMICS D6331, two Illumina replicates; Section~\ref{sec:mock}),
with a matched Kraken\,2$+$Bracken baseline, shows all methods dropping to
$r\approx0.41$--$0.58$ with overlapping uncertainty ($n{=}11$ in-set genera) and
Kraken\,2 abstaining on $\sim$56\% of reads---a substantial closed-set-to-real
gap. Replicate-to-replicate agreement is tight, but the community is still a
single standard (11 in-set genera, 135\,bp reads); a fuller evaluation across
multiple communities, matched read lengths, and additional baselines
(e.g.\ Centrifuge) remains the priority follow-up.

\paragraph{Baseline breadth.}
We compare against Kraken\,2 both as raw report counts and with Bracken
\citep{Lu2017} abundance re-estimation (Section~\ref{sec:sample}); Bracken is the
appropriate abundance baseline and is strong ($r=0.997$ on the simulated
in-database pool). A fuller benchmark should also add Centrifuge
\citep{Kim2016} or other classifiers so that the neural models are measured
against the broadest set of production pipelines.

\paragraph{Single read length and evaluation asymmetry.}
Simulated experiments use 150\,bp reads; the real mock probe uses median
135\,bp reads from the same Illumina run. Generalization to other read lengths
(75/250\,bp) is untested. We also note a metric asymmetry in the current
results---NT-v2 numbers use reverse-complement test-time augmentation while
MetaTransformer numbers are forward-only---which slightly understates the
MetaTransformer models (RC-TTA adds $\sim$0.7 pp); harmonizing this is
straightforward and does not affect any qualitative conclusion.

\paragraph{Outlook.}
The consistent thread is that input representation, not model scale, is the
binding constraint for short-read taxonomy with foundation models. The most
promising direction is a microbial foundation model pre-trained with
overlapping long-$k$ tokenization and an efficient large-vocabulary embedding,
evaluated from the outset with the leakage control, out-of-genome splits, and
dual read/sample metrics used here.
